% GENERATED FILE. Edit canonical lessons, then run npm run generate:books.
% canonical-source-hash: fnv1a64:39c9b6dc96edde97
% canonical-lessons: KA-C383-sandeha, KA-C383-tirucu, KA-C383-nirnayisu, KA-C383-asadhya, KA-C383-allagale, KA-R383-review-376-383

\chapter{Doubts and Judgements}
\label{ch:ka-doubts-and-judgements}
\hlchaptermodality{383}

\begin{chapteropening}
\textbf{By the end of this chapter:} \emph{I can say a misgiving, to twist, to determine, impossible and to deny.}

Complete the last lesson of chapter 383: I can say a misgiving, to twist, to determine, impossible and to deny.
\end{chapteropening}

\section[\texorpdfstring{\kn{ಸಂದೇಹ} (sandēha)}{ಸಂದೇಹ (sandēha)}]{\kn{ಸಂದೇಹ} (sandēha) — a misgiving; a doubt}

\label{lesson:KA-C383-sandeha}



\begin{quote}
Before the new one: say the Kannada for existence, then the Kannada for to realise.
\end{quote}

\subsection*{You'll want to know: \kn{ಸಂದೇಹ}}
\textbf{\kn{ಸಂದೇಹ}} — \emph{sandēha} — \textquotedblleft{}a misgiving; a doubt\textquotedblright{}.

\begin{quote}
\kn{ಅವನ} \kn{ಕಥೆಯ} \kn{ಬಗ್ಗೆ} \kn{ನನಗೆ} \kn{ಸಂದೇಹ} \kn{ಇದೆ}.
\end{quote}

\begin{quote}
\emph{avana katheya bagge nanage sandēha ide.}
\end{quote}

\begin{quote}
I have doubts about his story.
\end{quote}

Who doubts takes -\kn{ಗೆ}: \kn{ನನಗೆ} \kn{ಸಂದೇಹ} \kn{ಇದೆ}.

\begin{cousinweb}
\textbf{\kn{ಸಂದೇಹ}} (\emph{sandēha}) is borrowed from Sanskrit \emph{saṃdeha}, \textquotedblleft{}doubt, uncertainty\textquotedblright{} (Monier-Williams).
\end{cousinweb}

\begin{grammarlens}[title={What you've built}]
One more word for giving reasons and opinions.
\end{grammarlens}

\subsection*{Your turn}
\begin{itemize}
\raggedright
  \item \emph{Say it:} \emph{sandēha}
  \item \emph{Say it:} \emph{sandēha}, once more
  \item \emph{Say it:} \emph{avana katheya bagge nanage sandēha ide.} — I have doubts about his story.
  \item \emph{Recall:} say the Kannada for existence, then the Kannada for to realise, then say \emph{sandēha} again
\end{itemize}

\begin{tcolorbox}[breakable,colback=teal!4,colframe=teal!35!black,title={Before you move on}]
What does \kn{ಸಂದೇಹ} mean? (\textquotedblleft{}A misgiving; a doubt\textquotedblright{}.) Say it once more.
\end{tcolorbox}
\section[\texorpdfstring{\kn{ತಿರುಚು} (tirucu)}{ತಿರುಚು (tirucu)}]{\kn{ತಿರುಚು} (tirucu) — to twist; to distort}

\label{lesson:KA-C383-tirucu}



\begin{quote}
Before the new one: say the Kannada for to realise, then the Kannada for a misgiving.
\end{quote}

\subsection*{You'll want to know: \kn{ತಿರುಚು}}
\textbf{\kn{ತಿರುಚು}} — \emph{tirucu} — \textquotedblleft{}to twist; to distort\textquotedblright{}.

\begin{quote}
\kn{ಆ} \kn{ಪತ್ರಿಕೆ} \kn{ಅವರ} \kn{ಮಾತನ್ನು} \kn{ತಿರುಚಿತು}.
\end{quote}

\begin{quote}
\emph{ā patrike avara mātannu tirucitu.}
\end{quote}

\begin{quote}
That newspaper twisted their words.
\end{quote}

Past (I …): \kn{ತಿರುಚಿದೆ}, \emph{tirucide}.

Of facts, words and limbs: \kn{ಕೈ} \kn{ತಿರುಚು}, to twist an arm.

\begin{cousinweb}
\textbf{\kn{ತಿರುಚು}} (\emph{tirucu}) is a Kannada verb whose origin is not settled: DEDR does not list it.
\end{cousinweb}

\begin{grammarlens}[title={What you've built}]
One more word for giving reasons and opinions.
\end{grammarlens}

\subsection*{Your turn}
\begin{itemize}
\raggedright
  \item \emph{Say it:} \emph{tirucu}
  \item \emph{Say it:} \emph{tirucu}, once more
  \item \emph{Say it:} \emph{ā patrike avara mātannu tirucitu.} — That newspaper twisted their words.
  \item \emph{Recall:} say the Kannada for to realise, then the Kannada for a misgiving, then say \emph{tirucu} again
\end{itemize}

\begin{tcolorbox}[breakable,colback=teal!4,colframe=teal!35!black,title={Before you move on}]
What does \kn{ತಿರುಚು} mean? (\textquotedblleft{}To twist; to distort\textquotedblright{}.) Say it once more.
\end{tcolorbox}
\section[\texorpdfstring{\kn{ನಿರ್ಣಯಿಸು} (nirṇayisu)}{ನಿರ್ಣಯಿಸು (nirṇayisu)}]{\kn{ನಿರ್ಣಯಿಸು} (nirṇayisu) — to determine; to judge}

\label{lesson:KA-C383-nirnayisu}



\begin{quote}
Before the new one: say the Kannada for a misgiving, then the Kannada for to twist.
\end{quote}

\subsection*{You'll want to know: \kn{ನಿರ್ಣಯಿಸು}}
\textbf{\kn{ನಿರ್ಣಯಿಸು}} — \emph{nirṇayisu} — \textquotedblleft{}to determine; to judge\textquotedblright{}.

\begin{quote}
\kn{ಪುರಾವೆಗಳನ್ನು} \kn{ನೋಡಿ} \kn{ಅವರು} \kn{ಸತ್ಯವನ್ನು} \kn{ನಿರ್ಣಯಿಸಿದರು}.
\end{quote}

\begin{quote}
\emph{purāvegaḷannu nōḍi avaru satyavannu nirṇayisidaru.}
\end{quote}

\begin{quote}
Looking at the evidence, they determined the truth.
\end{quote}

Past (I …): \kn{ನಿರ್ಣಯಿಸಿದೆ}, \emph{nirṇayiside}.

\kn{ನಿರ್ಣಯ}, a decision or a resolution.

\begin{cousinweb}
\textbf{\kn{ನಿರ್ಣಯಿಸು}} (\emph{nirṇayisu}) is built on Sanskrit \emph{nirṇaya}, \textquotedblleft{}complete ascertainment, decision\textquotedblright{} (Monier-Williams), with the Kannada verb ending -\emph{isu}.
\end{cousinweb}

\begin{grammarlens}[title={What you've built}]
One more word for giving reasons and opinions.
\end{grammarlens}

\subsection*{Your turn}
\begin{itemize}
\raggedright
  \item \emph{Say it:} \emph{nirṇayisu}
  \item \emph{Say it:} \emph{nirṇayisu}, once more
  \item \emph{Say it:} \emph{purāvegaḷannu nōḍi avaru satyavannu nirṇayisidaru.} — Looking at the evidence, they determined the truth.
  \item \emph{Recall:} say the Kannada for a misgiving, then the Kannada for to twist, then say \emph{nirṇayisu} again
\end{itemize}

\begin{tcolorbox}[breakable,colback=teal!4,colframe=teal!35!black,title={Before you move on}]
What does \kn{ನಿರ್ಣಯಿಸು} mean? (\textquotedblleft{}To determine; to judge\textquotedblright{}.) Say it once more.
\end{tcolorbox}
\section[\texorpdfstring{\kn{ಅಸಾಧ್ಯ} (asādhya)}{ಅಸಾಧ್ಯ (asādhya)}]{\kn{ಅಸಾಧ್ಯ} (asādhya) — impossible}

\label{lesson:KA-C383-asadhya}



\begin{quote}
Before the new one: say the Kannada for to twist, then the Kannada for to determine.
\end{quote}

\subsection*{You'll want to know: \kn{ಅಸಾಧ್ಯ}}
\textbf{\kn{ಅಸಾಧ್ಯ}} — \emph{asādhya} — \textquotedblleft{}impossible\textquotedblright{}.

\begin{quote}
\kn{ಒಂದೇ} \kn{ದಿನದಲ್ಲಿ} \kn{ಈ} \kn{ಕೆಲಸ} \kn{ಮುಗಿಸುವುದು} \kn{ಅಸಾಧ್ಯ}.
\end{quote}

\begin{quote}
\emph{ondē dinadalli ī kelasa mugisuvudu asādhya.}
\end{quote}

\begin{quote}
Finishing this job in a single day is impossible.
\end{quote}

\kn{ಅ}- (not) + \kn{ಸಾಧ್ಯ} (possible).

\begin{cousinweb}
\textbf{\kn{ಅಸಾಧ್ಯ}} (\emph{asādhya}) is borrowed from Sanskrit \emph{asādhya}, \textquotedblleft{}not to be accomplished\textquotedblright{}, from \emph{a-}, \textquotedblleft{}not\textquotedblright{}, and \emph{sādhya}, \textquotedblleft{}to be done\textquotedblright{} (Monier-Williams).
\end{cousinweb}

\begin{grammarlens}[title={What you've built}]
One more word for giving reasons and opinions.
\end{grammarlens}

\subsection*{Your turn}
\begin{itemize}
\raggedright
  \item \emph{Say it:} \emph{asādhya}
  \item \emph{Say it:} \emph{asādhya}, once more
  \item \emph{Say it:} \emph{ondē dinadalli ī kelasa mugisuvudu asādhya.} — Finishing this job in a single day is impossible.
  \item \emph{Recall:} say the Kannada for to twist, then the Kannada for to determine, then say \emph{asādhya} again
\end{itemize}

\begin{tcolorbox}[breakable,colback=teal!4,colframe=teal!35!black,title={Before you move on}]
What does \kn{ಅಸಾಧ್ಯ} mean? (\textquotedblleft{}Impossible\textquotedblright{}.) Say it once more.
\end{tcolorbox}
\section[\texorpdfstring{\kn{ಅಲ್ಲಗಳೆ} (allagaḷe)}{ಅಲ್ಲಗಳೆ (allagaḷe)}]{\kn{ಅಲ್ಲಗಳೆ} (allagaḷe) — to deny; to refute}

\label{lesson:KA-C383-allagale}



\begin{quote}
Before the new one: say the Kannada for to determine, then the Kannada for impossible.
\end{quote}

\subsection*{You'll want to know: \kn{ಅಲ್ಲಗಳೆ}}
\textbf{\kn{ಅಲ್ಲಗಳೆ}} — \emph{allagaḷe} — \textquotedblleft{}to deny; to refute\textquotedblright{}.

\begin{quote}
\kn{ಅವನು} \kn{ಎಲ್ಲಾ} \kn{ಆರೋಪಗಳನ್ನು} \kn{ಅಲ್ಲಗಳೆದನು}.
\end{quote}

\begin{quote}
\emph{avanu ellā ārōpagaḷannu allagaḷedanu.}
\end{quote}

\begin{quote}
He denied all the accusations.
\end{quote}

Past (I …): \kn{ಅಲ್ಲಗಳೆದೆ}, \emph{allagaḷede}.

What you deny takes -\kn{ಅನ್ನು}.

\begin{cousinweb}
\textbf{\kn{ಅಲ್ಲಗಳೆ}} (\emph{allagaḷe}) is built from two parts: native \emph{alla}, \textquotedblleft{}it is not so\textquotedblright{} (DEDR 234), and \emph{kaḷe}, \textquotedblleft{}to cast off\textquotedblright{}: to throw a claim out as untrue.
\end{cousinweb}

\begin{grammarlens}[title={What you've built}]
One more word for giving reasons and opinions.
\end{grammarlens}

\subsection*{Your turn}
\begin{itemize}
\raggedright
  \item \emph{Say it:} \emph{allagaḷe}
  \item \emph{Say it:} \emph{allagaḷe}, once more
  \item \emph{Say it:} \emph{avanu ellā ārōpagaḷannu allagaḷedanu.} — He denied all the accusations.
  \item \emph{Recall:} say the Kannada for to determine, then the Kannada for impossible, then say \emph{allagaḷe} again
\end{itemize}

\begin{tcolorbox}[breakable,colback=teal!4,colframe=teal!35!black,title={Before you move on}]
What does \kn{ಅಲ್ಲಗಳೆ} mean? (\textquotedblleft{}To deny; to refute\textquotedblright{}.) Say it once more.
\end{tcolorbox}
\section[(dialogue)]{Reasons, proof and points of view}

\label{lesson:KA-R383-review-376-383}



\begin{quote}
Say the Kannada for proof and to prove.
\end{quote}

\subsection*{Your turn}
Say these aloud:

\begin{itemize}
\raggedright
  \item proof, to prove, a possibility, importance, to base on
  \item logic, to justify, a controversy, to object, reasonable
  \item a point of view, a stance, neutral, bias, to advocate
  \item to consent, a difference of opinion, a consensus, to oppose, mutual
  \item to persuade, to motivate, to drive home, an inducement, to attract
  \item to consider, to suppose, good sense, to reflect on, intellect
  \item a fact, reality, an illusion, existence, to realise
  \item a misgiving, to twist, to determine, impossible, to deny
\end{itemize}

\begin{tcolorbox}[breakable,colback=teal!4,colframe=teal!35!black,title={Before you move on}]
Proof? (\textbf{\kn{ಪುರಾವೆ}}, \emph{purāve}.) To prove? (\textbf{\kn{ಸಾಬೀತುಪಡಿಸು}}, \emph{sābītupaḍisu}.) A possibility? (\textbf{\kn{ಸಾಧ್ಯತೆ}}, \emph{sādhyate}.) Importance? (\textbf{\kn{ಮಹತ್ವ}}, \emph{mahatva}.) To base on? (\textbf{\kn{ಆಧರಿಸು}}, \emph{ādharisu}.) Logic? (\textbf{\kn{ತರ್ಕ}}, \emph{tarka}.) To justify? (\textbf{\kn{ಸಮರ್ಥಿಸು}}, \emph{samarthisu}.) A controversy? (\textbf{\kn{ವಿವಾದ}}, \emph{vivāda}.) To object? (\textbf{\kn{ಆಕ್ಷೇಪಿಸು}}, \emph{ākṣēpisu}.) Reasonable? (\textbf{\kn{ಸಮಂಜಸ}}, \emph{samañjasa}.) A point of view? (\textbf{\kn{ದೃಷ್ಟಿಕೋನ}}, \emph{dṛṣṭikōna}.) A stance? (\textbf{\kn{ನಿಲುವು}}, \emph{niluvu}.) Neutral? (\textbf{\kn{ತಟಸ್ಥ}}, \emph{taṭastha}.) Bias? (\textbf{\kn{ಪಕ್ಷಪಾತ}}, \emph{pakṣapāta}.) To advocate? (\textbf{\kn{ಪ್ರತಿಪಾದಿಸು}}, \emph{pratipādisu}.) To consent? (\textbf{\kn{ಸಮ್ಮತಿಸು}}, \emph{sammatisu}.) A difference of opinion? (\textbf{\kn{ಭಿನ್ನಾಭಿಪ್ರಾಯ}}, \emph{bhinnābhiprāya}.) A consensus? (\textbf{\kn{ಒಮ್ಮತ}}, \emph{ommata}.) To oppose? (\textbf{\kn{ವಿರೋಧಿಸು}}, \emph{virōdhisu}.) Mutual? (\textbf{\kn{ಪರಸ್ಪರ}}, \emph{paraspara}.) To persuade? (\textbf{\kn{ಮನವೊಲಿಸು}}, \emph{manavolisu}.) To motivate? (\textbf{\kn{ಪ್ರೇರೇಪಿಸು}}, \emph{prērēpisu}.) To drive home? (\textbf{\kn{ಮನದಟ್ಟು} \kn{ಮಾಡು}}, \emph{manadaṭṭu māḍu}.) An inducement? (\textbf{\kn{ಆಮಿಷ}}, \emph{āmiṣa}.) To attract? (\textbf{\kn{ಆಕರ್ಷಿಸು}}, \emph{ākarṣisu}.) To consider? (\textbf{\kn{ಪರಿಗಣಿಸು}}, \emph{parigaṇisu}.) To suppose? (\textbf{\kn{ಭಾವಿಸು}}, \emph{bhāvisu}.) Good sense? (\textbf{\kn{ವಿವೇಕ}}, \emph{vivēka}.) To reflect on? (\textbf{\kn{ಮನನ} \kn{ಮಾಡು}}, \emph{manana māḍu}.) Intellect? (\textbf{\kn{ಬುದ್ಧಿ}}, \emph{buddhi}.) A fact? (\textbf{\kn{ಸಂಗತಿ}}, \emph{saṅgati}.) Reality? (\textbf{\kn{ವಾಸ್ತವ}}, \emph{vāstava}.) An illusion? (\textbf{\kn{ಭ್ರಮೆ}}, \emph{bhrame}.) Existence? (\textbf{\kn{ಅಸ್ತಿತ್ವ}}, \emph{astitva}.) To realise? (\textbf{\kn{ಮನಗಾಣು}}, \emph{managāṇu}.) A misgiving? (\textbf{\kn{ಸಂದೇಹ}}, \emph{sandēha}.) To twist? (\textbf{\kn{ತಿರುಚು}}, \emph{tirucu}.) To determine? (\textbf{\kn{ನಿರ್ಣಯಿಸು}}, \emph{nirṇayisu}.) Impossible? (\textbf{\kn{ಅಸಾಧ್ಯ}}, \emph{asādhya}.) To deny? (\textbf{\kn{ಅಲ್ಲಗಳೆ}}, \emph{allagaḷe}.)
\end{tcolorbox}
